\section*{NeurIPS Paper Checklist}

\begin{enumerate}

\item {\bf Claims}
    \item[] Question: Do the main claims made in the abstract and introduction accurately reflect the paper's contributions and scope?
    \item[] Answer: \answerYes{}
    \item[] Justification: The abstract and Section~\ref{sec:intro} list four contributions (taxonomy, injection-detection validation, benchmark with empirical results, released artifacts) that map directly to Sections~\ref{sec:taxonomy}, \ref{sec:injection}, \ref{sec:geo}, and the data-release statement in Section~\ref{sec:conclusion}.

\item {\bf Limitations}
    \item[] Question: Does the paper discuss the limitations of the work performed by the authors?
    \item[] Answer: \answerYes{}
    \item[] Justification: Section~\ref{sec:conclusion}.1 explicitly enumerates limitations: the OOF-prediction requirement, single feature substrate (ACS only), taxonomy coverage (six modes is a v1 catalogue with two additional modes deferred), threshold calibration history, and no operational deployment claims.

\item {\bf Theory assumptions and proofs}
    \item[] Question: For each theoretical result, does the paper provide the full set of assumptions and a complete (and correct) proof?
    \item[] Answer: \answerNA{}
    \item[] Justification: This paper presents an empirical taxonomy and benchmark; it does not claim novel theorems. The $R/S_{\text{sup}}/N$ decomposition is taken as a structural definition from~\citep{martin2026rsct}; the simplex constraint $R{+}S_{\text{sup}}{+}N{=}1$ is enforced by construction.

\item {\bf Experimental result reproducibility}
    \item[] Question: Does the paper fully disclose all the information needed to reproduce the main experimental results?
    \item[] Answer: \answerYes{}
    \item[] Justification: Section~\ref{sec:geo}.1 specifies the 27 ZCTA tasks, the 33 ACS features, and the three model families with their hyperparameters (detailed in Appendix~\ref{app:solver_families}). Appendix~\ref{app:nceiling_method} provides bootstrap details, seed-collapse procedure, and synthetic-validation construction. Appendix~\ref{app:compute} reports compute environment. The HF dataset~\citep{geocert2026} provides out-of-fold predictions for all models.

\item {\bf Open access to data and code}
    \item[] Question: Does the paper provide open access to the data and code, with sufficient instructions to faithfully reproduce the main experimental results?
    \item[] Answer: \answerYes{}
    \item[] Justification: The GeoRSCT benchmark, out-of-fold predictions for all geospatial families, the certificate library, the taxonomy specification, and the injection-detection harness are released as GeoCert v23.001 on Hugging Face~\citep{geocert2026} under CC-BY-4.0. The dataset card provides loading examples and reproduction instructions.

\item {\bf Experimental setting/details}
    \item[] Question: Does the paper specify all the training and test details necessary to understand the results?
    \item[] Answer: \answerYes{}
    \item[] Justification: Data splits (county-holdout interpolation, state-holdout extrapolation), hyperparameters (GBDT: 200 trees, learning rate 0.05, max depth 4; gate thresholds $\kappa_{\text{base}}{=}0.5$, $\lambda{=}0.4$), aggregation rules (median over $R$-labeled samples for $\kappa$ and $\alpha$), and label construction (residual terciles) are reported in Section~\ref{sec:geo} and Appendix~\ref{app:solver_families}.

\item {\bf Experiment statistical significance}
    \item[] Question: Does the paper report error bars suitably and correctly defined or other appropriate information about the statistical significance of the experiments?
    \item[] Answer: \answerYes{}
    \item[] Justification: N-ceiling estimates use state-level block bootstrap with $B{=}1{,}000$ resamples to construct $[2.5\text{th}, 97.5\text{th}]$ percentile CIs (Section~\ref{sec:nceiling}, Appendix~\ref{app:nceiling_method}); these CIs are 15--40\% wider than naive bootstrap to account for spatial autocorrelation. Injection-detection results are reported across eight independent seeds (Section~\ref{sec:injection}).

\item {\bf Experiments compute resources}
    \item[] Question: Does the paper provide sufficient information on the computer resources needed to reproduce the experiments?
    \item[] Answer: \answerYes{}
    \item[] Justification: Appendix~\ref{app:compute} reports compute requirements: \texttt{ml.m5.4xlarge} (16 vCPU, 64 GB RAM) instances on AWS SageMaker, with wall-clock times. No GPUs were required. The injection-detection harness runs in seconds on commodity hardware.

\item {\bf Code of ethics}
    \item[] Question: Does the research conducted in the paper conform, in every respect, with the NeurIPS Code of Ethics?
    \item[] Answer: \answerYes{}
    \item[] Justification: The research uses only publicly available, aggregated data (CDC PLACES at ZCTA level, ACS census variables) and does not involve human subjects, identifiable personal data, or restricted resources.

\item {\bf Broader impacts}
    \item[] Question: Does the paper discuss both potential positive societal impacts and negative societal impacts of the work performed?
    \item[] Answer: \answerYes{}
    \item[] Justification: Section~\ref{sec:conclusion}.2 discusses both. Positive: the protocol surfaces structural properties of evaluation, supporting more careful deployment decisions. Negative: deployment decisions issued by the audit could be applied uncritically. We emphasize the protocol is diagnostic and should inform, not replace, human judgment.

\item {\bf Safeguards}
    \item[] Question: Does the paper describe safeguards that have been put in place for responsible release of data or models that have a high risk for misuse?
    \item[] Answer: \answerNA{}
    \item[] Justification: The released artifacts---the GeoRSCT benchmark (already-public CDC/ACS data at aggregated ZCTA level), prediction tables, the certificate library, the taxonomy specification, and the injection-detection harness---do not pose high misuse risks. There are no pre-trained generative models, no scraped or sensitive data, and no model weights for systems that could be repurposed for harmful generation.

\item {\bf Licenses for existing assets}
    \item[] Question: Are the creators or original owners of assets used in the paper properly credited and are the license and terms of use explicitly mentioned and properly respected?
    \item[] Answer: \answerYes{}
    \item[] Justification: CDC PLACES and ACS data are public-domain US government datasets. The $R/S_{\text{sup}}/N$ decomposition is cited to~\citep{martin2026rsct}. PDFM~\citep{agarwal2025pdfm} and SatCLIP~\citep{klemmer2023satclip} are cited where their results are discussed.

\item {\bf New assets}
    \item[] Question: Are new assets introduced in the paper well documented and is the documentation provided alongside the assets?
    \item[] Answer: \answerYes{}
    \item[] Justification: The GeoCert v23.001 release on Hugging Face~\citep{geocert2026} includes a dataset card documenting all 27 tasks, 33 ACS features, three evaluation protocols, simplified ZCTA geometries, schema, provenance, and a CC-BY-4.0 license. The taxonomy specification is documented both in the paper (Section~\ref{sec:taxonomy}, Table~\ref{tab:taxonomy}) and in the repository.

\item {\bf Crowdsourcing and research with human subjects}
    \item[] Question: For crowdsourcing experiments and research with human subjects, does the paper include the full text of instructions given to participants and screenshots, if applicable, as well as details about compensation?
    \item[] Answer: \answerNA{}
    \item[] Justification: The research involves no crowdsourcing and no human subjects.

\item {\bf Institutional review board (IRB) approvals or equivalent for research with human subjects}
    \item[] Question: Does the paper describe potential risks incurred by study participants, whether such risks were disclosed to the subjects, and whether IRB approvals were obtained?
    \item[] Answer: \answerNA{}
    \item[] Justification: The research does not involve human subjects. Aggregated population health prevalence data (CDC PLACES) and census aggregates (ACS) at the ZCTA level cannot be used to identify individuals.

\item {\bf Declaration of LLM usage}
    \item[] Question: Does the paper describe the usage of LLMs if it is an important, original, or non-standard component of the core methods in this research?
    \item[] Answer: \answerNA{}
    \item[] Justification: LLMs do not appear as a core methodological component. The certificate library, the failure taxonomy, the injection-detection harness, and the N-ceiling estimator do not involve LLM components.

\end{enumerate}
